\section{16 августа 1914 г.}

\lp{20260604_194517507}

Дорогой мой. Какая радость была для нас Ваша телеграмма. Мы измучались, думая,
что вас тоже, как и многих, арестовали. И Коля, и я спать не могли и даже
боялись заговаривать о вас. Волновались ужасно. А когда тут стали возвращаться
русские, мы все пытались узнать не видал ли кто вас. Тут ежедневно печатались
списки русских за границей. Пришли обе ваши телеграммы, и два дня тому назад
письмо папаше, а

\lp{20260604_194539803}

\noindent
сейчас письмо мне и открытка Коле. Первая телеграмма пришла шесть дней назад, а
в Кузнецкий переулок -- два дня. Я сейчас же ответила вам на первую телеграмму,
которая очевидно шла очень долго (на ней не было числа отправления). Писем уже
не посылали, т.к. сказали, что не отправляют. Сейчас вероятно дойдет. Ваше
письмо помечено 8 августа, но не знаю какого стиля; если нового, то значит шло
страшно долго. Сегодня мы переехали в город. Жили на даче с грехом пополам т.к.
все время ездили в город.

Напишите, ради Бога, про себя! Ничего не знаю. Неужели Ядвига не с вами? Почему
вы об этои ни слова?

\lp{20260604_194539803}

Рахманинова тоже освободили.

За это время мы так много пережили, и так все свалилось неожиданно. Колю
призывали \rem{в армию}, и его особенно потрясло, что он должен идти нижним
чином и попасть в столь непривычную обстановку. Но, когда мы узнали, что
призываны Гольденвейзер, Сабанеев, Петя Страхов и вообще все, все, то стали
только хлопотать о том, чтобы Коле попасть куда-нибудь вместе с кем-нибудь из
них. Воинский начальник оказался сущим ангелом и вероятно устроил бы так, что
все учителя, профессора, неспособные к строевой службе, заняли бы канцелярские
места и таким образом могли бы действительно принести какую-нибудь пользу.
Неделю тянулась канитель с участком, казармами, канцелярией

\lp{20260604_194645876}

\noindent
института с утра до ночи, и наконец, когда было медицинское
освидетельствование, Коля пошел без зубов и заявил, что у него не хватает их
очень много, и его освободили. Коля с ним \rem{с военным начальником} всюду
ездил, и это Колю очень поддержало. Потом оказалось что Зяське тоже надо
явиться, но его то положение очень приятное, т.к. он какой-то кондуктор в Моск.
чертежной или упаковочной, не знаю, и ничего не делает. Их так много, что его,
кажется, отпустят. Он очень поправился за лето и этим обстоятельством (службой)
совсем не расстроен. Наоборот у Коли, например, такое чувство, что он даже не
может ничего не делать и хочет чем-нибудь помочь. Мы встречали

\lp{20260604_194645876}

\noindent
раненных. Не стану лучше говорить о впечатлении. Все, все, кто только может
делают у себя лазареты. У Марго тоже будет. Коля будет помогать вероятно в
распределении больных, кого куда. Мы шьем белье. Провожали солдат, наделяли чем
могли. Вообще мы больше ничем не живем сейчас. Сабурова отправляют в Польшу.
Коля должен явиться через две недели. Его отставка потеряла значение благодаря
какому-то новому закону. О нашем папе мы очень волновались, т.к. он больной,
всего боится и остался там без денег. И

\lp{20260604_194701429}

\noindent
только вчера пришла телеграмма из Берлина что ждать денег. Мы думали, что в
этих толпах, которые рвались в Россию, он погибнет. Маня возвращалась почти 1/2
месяца и приехала голая. Теперь уже слава Богу проезд возможен, но это все-таки
трудно, и лучше тем, которые могут, оставаться там. О Мариетте ничего не знаю.
Где то Ильины? О папе мы вам телеграфировали, думали не можете ли что узнать.
Сейчас по Колиной открытке узнала, что ваше письмо ко мне и открытка отправлена
8-го нового стиля, значит это шло 20 дней! Буду писать каждый день, может быть
почта будет теперь исправнее. Сегодня еще раз напишу обо всем, но все советуют
по-русски. Ваша Анна.

Мамаша страшно радовалась известию о вас. Она тоже исстрадалась.
